% ===================================================================
%  TABELO N-RO 12 — La stacidomo. — La fervojoj. — La ŝipoj.
%  Traduction 2026 d'apres texto/io, controlee sur texto/fr et
%  texto/es. Consigne : voir texto/eo/10-tabelo-01.tex.
% ===================================================================

%%K t12-apar-1 apar
\VUsaut{4.0mm}
\VUtitre{15.0pt}{60}{TABELO N\textsuperscript{o} 12}
\VUsaut{2.4mm}
\VUpk{11.4pt}{40}{La stacidomo. --- La fervojoj. --- La ŝipoj.}
\VUsaut{3.4mm}
\textsc{Noto.} --- \textit{La horaroj uzataj en ĉiuj nacioj estas diversaj, ni
do uzas kiel ekzemplon la horojn de la} \VUgras{franca tabelo.}

%%K t12-01-1 p
1. --- Mi ĵus vojaĝis tra Germanio. Antaŭ la foriro, mi aĉetis
\VUgras{horlibron} \textsuperscript{(15)} por scii la foriran horon de la
trajnoj. Post konstati ke la plej oportuna trajno estas forironta el Parizo, je
naŭ horoj vespere, kaj alvenonta en Strasburgon je unu horo dek tri minutoj, mi
preparis mian vojaĝon, kaj en la sekva tago, mi igis konduki min al la
stacidomo.

%%K t12-02-1 p
2. --- Kiam mi eniris la grandan forirhalon, malantaŭ maljuna \VUgras{pastro}
\textsuperscript{(2)} kampara, kiu mane kunportis sian \VUgras{vojaĝsakon}
\textsuperscript{(3)}, multaj \VUgras{vojaĝontoj} \textsuperscript{(1)} estis jam
alvenintaj.

%%K t12-02-2 p
Ĉar mi volis sendi \VUgras{depeŝon (telegramon)} \textsuperscript{(5)}, mi igis
montri al mi de \VUgras{kontrolisto} \textsuperscript{(6)} \VUgras{la
telegrafoficejon} \textsuperscript{(4)}.

%%K t12-03-1 p
3. --- Antaŭ pasi en la \VUgras{atendosalonon} \textsuperscript{(7)} mi iris
provizi min per \VUgras{bileto} \textsuperscript{(9)} kun reveno.

%%K t12-04-1 p
4. --- Mi ne longe atendis ĉe la \VUgras{giĉeto} \textsuperscript{(8)}, ĉar
malmultaj personoj estis tie. \VUgras{Soldato} \textsuperscript{(10)}, kiu estis
foriranta forpermese, komplezis cedi al mi sian vicon, tiel ke mi havis
sufiĉan tempon por peti kelkajn informojn de la \VUgras{stacidomestro}
\textsuperscript{(13)}, kiu estis konversacianta kun la \VUgras{komisaro}
\textsuperscript{(12)} de la gvatado kaj kun la \VUgras{ĝendarmo}
\textsuperscript{(11)} deĵoranta.

%%K t12-tit-1 sub
\VUpk{10.2pt}{40}{Bagaĝenskribado.}
\VUsaut{3.0mm}

%%K t12-05-1 p
5. --- Aĉetinte ĵurnalon ĉe la \VUgras{biblioteko} \textsuperscript{(14)}, mi
igis pesi miajn \VUgras{bagaĝojn} \textsuperscript{(17)} kiujn \VUgras{portisto}
\textsuperscript{(16)} ĵus alkondukis per \VUgras{ĉaro pakaĵa}
\textsuperscript{(19)}; \VUgras{la oficisto} \textsuperscript{(22)} komisiita pri
la \VUgras{pesaparato} \textsuperscript{(21)} sciigis al mi, ke mia kofro pezas
tridek kvin kilogramojn; tio faris \VUgras{superpezon} de kvin kilogramoj; pro
kiu mi ŝuldas suplementan pagon, ĉe la \VUgras{giĉeto} de la bagaĝoj
\textsuperscript{(23)}.

%%K t12-06-1 p
6. --- Ĉar mi estis tro frua, mi havis liberan tempon por ekzameni la
cirkuladon de la vojaĝantoj en la stacidomo. Unuj deponis aŭ reprenis siajn
negravajn bagaĝojn ĉe la \VUgras{(bagaĝ)deponejo} \textsuperscript{(25)}; aliaj
konsultis la \VUgras{horarojn} \textsuperscript{(26)}. La alvenantaj vojaĝantoj
rapidis al la \VUgras{elirejo} \textsuperscript{(32)} kun sia \VUgras{valizo}
\textsuperscript{(24)} en la mano. Kelkaj atendis en la
\VUgras{(bagaĝ)distribuejo} \textsuperscript{(27)} kaj prezentis sian
\VUgras{bileton} \textsuperscript{(29)} por repreni siajn kofrojn,
\VUgras{kestojn} \textsuperscript{(28)} aŭ \VUgras{korbojn}
\textsuperscript{(30)}.

%%K t12-07-1 p
7. --- \VUgras{La oficistoj de la akcizo} \textsuperscript{(33)} zorge esploris
la pakojn por konstati ĉu oni havas ion deklarindan.

%%K t12-08-1 p
8. --- Iuj vojaĝantoj iris al la \VUgras{bufedo} \textsuperscript{(34)}, kie
\VUgras{kelnero} \textsuperscript{(35)} fervoris pri servi ilin. Ili devas
rapidi, ĉar la \VUgras{trajno} \textsuperscript{(36)}, kiu estas ekspreso, ne
longe restos en la stacidomo.

%%K t12-09-1 p
9. --- Poste mi iris al la fortrotuaro kaj serĉis rektan \VUgras{vagonon}
\textsuperscript{(37)} por ne esti devigata ŝanĝi la vagonon dum la vojaĝo. Mi
supreniris en koridorveturilon kiu enhavas \VUgras{fakon}
\textsuperscript{(38)} por fumantoj kaj fakon rezervitan por « sinjorinoj
solaj ».

%%K t12-tit-2 sub
\VUpk{10.2pt}{40}{Ekiro dum la nokto.}
\VUsaut{3.0mm}

%%K t12-10-1 p
10. --- Post supreniri la \VUgras{piedapogilon} \textsuperscript{(40)}, fermi la
\VUgras{pordon} \textsuperscript{(39)}, volvi la \VUgras{kurtenon}
\textsuperscript{(41)} kaj loki miajn malgrandajn bagaĝojn sur la \VUgras{reton}
\textsuperscript{(43)}, mi sidis sur la \VUgras{benketo} \textsuperscript{(42)}.
Vidante la \VUgras{lampiston} \textsuperscript{(44)} ekbruligi la lampojn, mi
pensis, ke mi devos pasigi nokton en la vagono kaj mi vokis la oficiston
komisiitan pri la lupago de la \VUgras{kapkusenoj} \textsuperscript{(45)} por
peti unu el ili.

%%K t12-11-1 p
11. --- La trajnkonduktoro venis kontroli la biletojn. Mi demandis lin kial ni
ne jam foriris, ĉar estas la ĝusta horo. Li respondis al mi ke la
\VUgras{trajnestro} \textsuperscript{(46)} okupiĝas pri metigi biciklojn en la
\VUgras{furgonon} \textsuperscript{(47)}. Krome, oni ne jam ŝanĝis ĉiujn
piedvarmigilojn \textsuperscript{(48)}.

%%K t12-12-1 p
12. --- Sed ni estis baldaŭ forironta, ĉar la \VUgras{hejtisto}
\textsuperscript{(50)}, sur sia \VUgras{tendro} \textsuperscript{(49)} prenis
karbon por eksciti la fajron de la \VUgras{lokomotivo} \textsuperscript{(54)},
kaj \VUgras{la maŝinisto} \textsuperscript{(51)} atendis nur la signalon de la
\VUgras{stacidoma subestro} \textsuperscript{(52)}, kiu estis sur la
\VUgras{trotuaro} \textsuperscript{(53)}.

%%K t12-13-1 p
13. --- La \VUgras{kaldrono} \textsuperscript{(56)} de la lokomotivo obtuze
grumblis; la \VUgras{kameno} \textsuperscript{(55)} vomis torentojn da fumo
miksita kun \VUgras{vaporo} \textsuperscript{(60)}. Stridanta \VUgras{fajfo}
\textsuperscript{(59)} sonegis kaj la \VUgras{pistonoj} \textsuperscript{(58)}
ekmoviĝis, impulsante la \VUgras{radojn} \textsuperscript{(57)} de la peza
maŝino.

%%K t12-tit-3 sub
\VUsaut{3.0mm}
\VUpk{10.2pt}{40}{Foriro!}
\VUsaut{3.0mm}

%%K t12-14-1 p
14. --- La rapideco de la trajno, kuranta sur la \VUgras{reloj}
\textsuperscript{(64)} estis akcelata; la \VUgras{trotuaroj}
\textsuperscript{(65)} malaperis; ni preterpasis la \VUgras{necesejojn}
\textsuperscript{(66)}, kaj ankaŭ vagonaron garitan sur la alia \VUgras{vojo}
\textsuperscript{(63)}: \VUgras{dormvagonoj} \textsuperscript{(67)},
\VUgras{restoraciovagono} \textsuperscript{(68)}, \VUgras{poŝtvagono}
\textsuperscript{(69)}.

%%K t12-noto-1 noto
\VUnotes{143.2mm}{%
\textit{(*) Noto.} --- \textit{Kompreneble la vojaĝopriskribo estas laŭ la
antaŭmilita landlimo.}%
}

%%K t12-15-1 p
15. --- Poste, la \VUgras{stacidomo de la varoj} \textsuperscript{(71)} pasis
antaŭ niaj okuloj. Sub la hangaroj, oficistoj ŝarĝis la \VUgras{varojn}
\textsuperscript{(72)} ekspeditajn per malrapida trajno. \VUgras{Eskvadranoj}
\textsuperscript{(76)} apud la \VUgras{akvorezervujo} \textsuperscript{(73)}
manovris \VUgras{bestovagonojn} \textsuperscript{(75)} kaj \VUgras{platformajn
vagonojn} \textsuperscript{(79)} por formi \VUgras{varotrajnon}
\textsuperscript{(80)}.

%%K t12-16-1 p
16. --- Ni transiris la lastan \VUgras{disbranĉiĝon} \textsuperscript{(78)},
apud kiu staris la relŝanĝisto; ni preteriris la \VUgras{signaldiskon}
\textsuperscript{(86)} kaj la stacidomo malproksime malaperis.

%%K t12-17-1 p
17. --- Baldaŭ ni trairis \VUgras{feran ponton} \textsuperscript{(74)} kiu
transpasas la \VUgras{fluon} \textsuperscript{(81)}. Post trairi tiun ponton,
ni laŭiris la riveron, sur kiu \VUgras{remorkiloj} \textsuperscript{(82)}
potencaj tiris pezajn \VUgras{kargoboatojn} \textsuperscript{(83)}. Ni krome
vidis \VUgras{vojaĝŝipon} \textsuperscript{(84)}, kiam ĝi estis albordiĝanta al
\VUgras{pontono} \textsuperscript{(85)}.

%%K t12-18-1 p
18. --- Post rapidege preterpasi plurajn stacidomojn, ni trairis kelkajn
\VUgras{tunelojn} \textsuperscript{(87)} kaj lasis malantaŭ ni \VUgras{dometojn}
sennombrajn \textsuperscript{(88)} de \VUgras{barilgardistoj}. Lulate de la
skuado de la trajno, mi profunde ekdormis.

%%K t12-tit-4 sub
\VUsaut{3.0mm}
\VUpk{10.2pt}{40}{Deutsch-Avricourt-Stacidomo.}
\VUsaut{3.0mm}

%%K t12-19-1 p
19. --- Mi estis subite vekigata de la voĉo de oficisto krianta:
« Deutsch-Avricourt! \textsuperscript{(*)}. Ĉiuj malsupreniru! » Okazis la
inspekto de la dogano kaj ĉiuj tedaj formalaĵoj de la landlimo. Mi devis
konformigi mian poŝhorloĝon al la \VUgras{horloĝego} \textsuperscript{(89)} de
la stacidomo, kiu estis tro frua je kvindek kvin minutoj; apenaŭ vekigita, mi
malfacile distingis la \VUgras{grandan montrilon} \textsuperscript{(90)} de la
\VUgras{malgranda} \textsuperscript{(91)}; la \VUgras{ciferplato}
\textsuperscript{(92)} mem estis tute konfuza pro la matena nebuleto.

%%K t12-20-1 p
20. --- Dum la doganistoj faris sian inspekton, mi oscedante deĉifris la
\VUgras{afiŝojn} \textsuperscript{(93)} kiuj dekoras la stacidomajn murojn. Mi
tagmanĝis por pasigi la tempon, konsultis la mapon de la germana \VUgras{reto}
\textsuperscript{(94)} por vidi kiom da distanco ankoraŭ apartigas min de
Strasburgo, kaj mi reeniris mian vagonon, refortigita de la espero baldaŭ
atingi la celon de mia vojaĝo.
\VUsaut{6.0mm}
\VUfilet{22mm}
